\documentclass[11pt]{article}
\usepackage[a4paper,margin=26mm]{geometry}
\usepackage{amsmath,amssymb,amsthm}
\usepackage[T1]{fontenc}
\usepackage{booktabs}
\usepackage{array}
\usepackage{placeins}
\usepackage{xurl}
\usepackage{hyperref}
\hypersetup{colorlinks=true,linkcolor=blue,urlcolor=blue,citecolor=blue,
 pdftitle={A negative answer to a question of Carenini on almost independent sets in regular graphs},
 pdfauthor={Horace Dong}}
\newtheorem{theorem}{Theorem}
\newtheorem{lemma}[theorem]{Lemma}
\newtheorem{corollary}[theorem]{Corollary}
\newtheorem{proposition}[theorem]{Proposition}
\newtheorem*{thmA}{Theorem A}
\newtheorem*{thmB}{Theorem B}
\newtheorem*{thmBp}{Theorem B$'$}
\newtheorem*{propC}{Proposition C}
\newtheorem*{thmD}{Theorem D}
\newtheorem*{question}{Question}
\theoremstyle{remark}
\newtheorem{remark}[theorem]{Remark}
\newcommand{\Prb}{\mathbb P}
\newcommand{\Ex}{\mathbb E}
\newcommand{\Nle}[1]{N_{\le #1}}
\newcommand{\cart}{\mathbin{\square}}
\title{A negative answer to a question of Carenini\\ on almost independent sets in regular graphs}
\author{Horace Dong\\ \small Independent researcher\\
\small\texttt{maxwelldhx@gmail.com}\\
\small ORCID: \href{https://orcid.org/0009-0002-7554-7548}{0009-0002-7554-7548}}
\date{September 2026}
\begin{document}
\maketitle
\begin{abstract}
For a $d$-regular graph $G$ on $n$ vertices and $\gamma\ge0$, let $i_\gamma(G)$
be the number of vertex sets that span at most $\gamma dn$ edges. Carenini
asked whether, when $2d\mid n$, the disjoint union of $n/(2d)$ copies of
$K_{d,d}$ maximises $i_\gamma$ among $d$-regular graphs on $n$ vertices for
every $\gamma\ge0$. For $\gamma=0$ this is the theorem of Kahn and Zhao on
independent sets. We show that the answer is no. The smallest counterexample
(smallest $n$, and smallest $\gamma$ for that $n$) is
$(n,d,\gamma)=(6,3,1/18)$: the triangular prism has $28$ vertex sets
spanning at most one edge, and $K_{3,3}$ has $24$. More generally, for every
pair $(n,d)$ with $2d\mid n$, except the trivial cases $d=1$ and
$(n,d)=(4,2)$ in which only one $d$-regular graph exists, the answer is
negative for some $\gamma<1/2$; for $d\ge3$ the union of copies of $K_{d,d}$
is then even a minimiser, and every $d$-regular graph containing a triangle
beats it. For small $\gamma$ we give two families: a switched $K_{d,d}$ beats
$K_{d,d}$ on $2d$ vertices for every $d\ge3$, and a connected bipartite
$d$-regular graph beats two copies of $K_{d,d}$ on $4d$ vertices for every
$d\ge2$, so the answer is no even among bipartite graphs. For every fixed
$\gamma\in(1/8,1/2)$ and every $d\ge2$, a bipartite competitor wins for all
sufficiently large $n$. On the other hand, the union of copies of $K_{d,d}$ is
optimal at the exponential scale for every $\gamma$; the upper bound here is
Carenini's biclique reduction, which rests on a theorem of Sah, Sawhney,
Stoner and Zhao. The exact question remains open for fixed
$\gamma\in(0,1/8]$ as $n\to\infty$, and our counterexamples are consistent
with the weaker statement in which the union of copies of $K_{d,d}$ is
optimal up to a factor $1+o(1)$. The counterexample at $(6,3,1/18)$, the
switched family for every $d\ge3$ and three further instances are formally
verified in Lean~4. The counterexamples and proofs were found by AI coding
agents working under the author's direction; their role is described at the
end of the paper.
\end{abstract}

\section{Introduction}

\paragraph{Background.}
For a graph $G$ and a set $A\subseteq V(G)$ let $e_G(A)$ be the number of
edges of $G$ with both ends in $A$. Following Carenini~\cite{car}, for a
$d$-regular graph $G$ on $n$ vertices and a real number $\gamma\ge0$ put
\[
 i_\gamma(G)=\bigl|\{A\subseteq V(G):\ e_G(A)\le\gamma dn\}\bigr|,
\]
so that $i_0(G)$ is the number of independent sets of~$G$. Alon~\cite{alon}
proved that a $d$-regular graph on $n$ vertices has at most
$2^{(1/2+o(1))n}$ independent sets, where $o(1)\to0$ as $d\to\infty$.
Kahn~\cite{kahn} proved the sharp bound $i_0(G)\le(2^{d+1}-1)^{n/(2d)}$ for
bipartite $d$-regular graphs, and Zhao~\cite{zhao} extended it to all
$d$-regular graphs. Since $K_{d,d}$ has $2^{d+1}-1$ independent sets, this
says that when $2d\mid n$, the disjoint union $mK_{d,d}$ of $m=n/(2d)$ copies
of $K_{d,d}$ has the largest number of independent sets among $d$-regular
graphs on $n$ vertices.

In work on tolerant testing of independent sets, Seth~\cite{seth} bounded the
number of sets $J\subseteq V(G)$ whose edge density $e_G(J)/\binom{|J|}2$ is
at most $d/(kn)$, and asked for the correct order of this bound and whether
the disjoint union of $n/(2d)$ copies of $K_{d,d}$ is extremal.
Carenini~\cite{car} proved a robust form of the Kahn--Zhao theorem for the
count $i_\gamma$: there is an absolute constant $C$ such that
\[
 i_\gamma(G)\le2^{n/2}\exp\bigl\{C\gamma\log(e/\gamma)\,n+Cn/d\bigr\}
\]
for every $d$-regular graph $G$ on $n$ vertices with $d\ge2$ and every
$0<\gamma\le1/2$ \cite[Theorem~1.1]{car}, and both error terms are
needed~\cite[Proposition~1.2]{car}. As Carenini notes, Seth's density
condition implies $e_G(J)\le dn/(2k)$, so this answers the quantitative part
of Seth's question. The key step~\cite[Proposition~2.1]{car} is a biclique
reduction based on a theorem of Sah, Sawhney, Stoner and
Zhao~\cite[Corollary~1.15]{sssz} (Corollary~1.14 in version~2 of the arXiv
preprint, the numbering used in~\cite{car}): for $0<q\le1$, the sum of
$q^{e_G(A)}$ over
all $A\subseteq V(G)$ is at most the $n/(2d)$-th power of the sum of
$q^{e(A)}$ over all subsets $A$ of $V(K_{d,d})$. Carenini then asked for the
exact extremal statement, which we restate in our notation.

\begin{question}[{\cite[Question~1.3]{car}}]
Suppose that $2d\mid n$, and let $m=n/(2d)$. Is it true that for every
$\gamma\ge0$, every $d$-regular graph $G$ on $n$ vertices satisfies
$i_\gamma(G)\le i_\gamma(mK_{d,d})$?
\end{question}

For $\gamma=0$ this is the theorem of Kahn and Zhao. For $\gamma\ge1/2$ it
holds trivially, since every set spans at most $dn/2\le\gamma dn$ edges and so
$i_\gamma(G)=2^n$. We call a pair $(n,d)$ \emph{admissible} if $d\ge1$ and
$2d\mid n$. Seth's count differs from $i_\gamma$, because its density
condition depends on $|J|$; our results concern Question~1.3, and we make no
claim about Seth's count.

\paragraph{Results.}
For an integer $t$ let $\Nle t(G)$ be the number of sets $A\subseteq V(G)$
with $e_G(A)\le t$. Since $e_G(A)$ is an integer,
$i_\gamma(G)=\Nle{\lfloor\gamma dn\rfloor}(G)$. Hence an inequality
$\Nle t(G)>\Nle t(mK_{d,d})$ at an integer $t\ge0$ gives a negative answer to
Question~1.3 for every $\gamma\in[t/(dn),(t+1)/(dn))$. We write $T(G)$ for
the number of triangles of $G$, $C_k$ for the cycle of length $k$, and
$G\cart H$ for the Cartesian product.

\begin{thmA}
Let $d\ge2$, let $n$ be a multiple of $2d$ with $(n,d)\ne(4,2)$, let
$m=n/(2d)$, and put
\[
 t^*=\frac{dn}2-3d+1,\qquad
 \gamma^*=\frac{t^*}{dn}=\frac12-\frac{3d-1}{dn}.
\]
Then $t^*\ge1$, and the following hold.
\begin{itemize}
\item[(a)] If $d\ge3$, every $d$-regular graph $G$ on $n$ vertices satisfies
 $i_{\gamma^*}(G)-i_{\gamma^*}(mK_{d,d})\ge\frac75T(G)$, and
 $i_{\gamma^*}(G)=i_{\gamma^*}(mK_{d,d})$ if and only if $G$ is
 triangle-free. So $mK_{d,d}$ minimises $i_{\gamma^*}$ among $d$-regular graphs
 on $n$ vertices, and every such graph that contains a triangle has a larger
 value.
\item[(b)] If $d=2$, every $2$-regular graph $G\ne mC_4$ on $n$ vertices
 satisfies $i_{\gamma^*}(G)>i_{\gamma^*}(mC_4)$.
\end{itemize}
\end{thmA}

\begin{table}[t]\centering\small
\begin{tabular}{@{}lcccccccccc@{}}\toprule
$t$ & 0 & 1 & 2 & 3 & 4 & 5 & 6 & 7 & 8 & 9\\\midrule
$\Nle t(K_{3,3})$ & 15 & 24 & 42 & 48 & 57 & 57 & 63 & 63 & 63 & 64\\
$\Nle t(K_3\cart K_2)$ & 13 & 28 & 40 & 48 & 57 & 57 & 63 & 63 & 63 & 64\\
\bottomrule
\end{tabular}
\caption{The number of vertex sets spanning at most $t$ edges, for the two
cubic graphs on six vertices. Here $i_\gamma=\Nle t$ with
$t=\lfloor18\gamma\rfloor$.}\label{tab:six}
\end{table}

The excluded cases are trivial: for $d=1$ the only $1$-regular graph on $n$
vertices is $\frac n2K_2=\frac n2K_{1,1}$, and $C_4=K_{2,2}$ is the only
$2$-regular graph on four vertices. For every other admissible pair there is a
competitor. If $d\ge3$, the graph $(K_d\cart K_2)\cup(m-1)K_{d,d}$ is
$d$-regular and contains a triangle; if $d=2$ and $n\ge8$, the cycle $C_n$
differs from $mC_4$. So \emph{the answer to Question~1.3 is negative for every
admissible pair $(n,d)$ except the trivial ones.}

For $(n,d)=(6,3)$ we get $t^*=1$ and $\gamma^*=1/18$. The only cubic graphs on
six vertices are $K_{3,3}$ and the triangular prism $K_3\cart K_2$, and
Table~\ref{tab:six} lists their counts. The prism has $28$ sets spanning at
most one edge, and $K_{3,3}$ has $24$. So $K_{3,3}$ fails to be a maximiser
exactly for $\gamma\in[1/18,1/9)$. Since $n=6$ is the smallest order at which
Question~1.3 is not trivial, $(n,d,\gamma)=(6,3,1/18)$ is the smallest
counterexample, in the sense that it has the smallest $n$ and, for that $n$,
the smallest $\gamma$.

For small $\gamma$ we give two explicit families.

\begin{thmB}
Let $d\ge3$, let $K_{d,d}$ have sides $X=\{x_1,\dots,x_d\}$ and
$Y=\{y_1,\dots,y_d\}$, and let $S_d$ be obtained from $K_{d,d}$ by deleting
the edges $x_1y_1$ and $x_2y_2$ and adding the edges $x_1x_2$ and $y_1y_2$.
Then $S_d$ is $d$-regular, every set spanning at most one edge of $K_{d,d}$
spans at most one edge of $S_d$, and
\[
 \Nle1(S_d)=\Nle1(K_{d,d})+4d-8=2^{d+1}+d^2+4d-9.
\]
Consequently, for $n=2d$ and every $\gamma\in[1/(2d^2),1/d^2)$, $K_{d,d}$
does not maximise $i_\gamma$ among $d$-regular graphs on $n$ vertices.
\end{thmB}

The graph $S_3$ is the prism. Corollary~\ref{cor:copies} gives an explicit
criterion for $S_d\cup(m-1)K_{d,d}$ to beat $mK_{d,d}$.

\begin{thmBp}
Let $d\ge2$. Take two disjoint copies $(X_1,Y_1)$ and $(X_2,Y_2)$ of
$K_{d,d}$, choose $u_k\in X_k$ and $w_k\in Y_k$ for $k=1,2$, delete the edges
$u_1w_1$ and $u_2w_2$, and add the edges $u_1w_2$ and $u_2w_1$. The resulting
graph $H_d$ is $d$-regular, bipartite and connected, and
\[
 \Nle1(H_d)-\Nle1(2K_{d,d})=2(d-1)(2^d+d-3)>0,\qquad
 i_0(H_d)-i_0(2K_{d,d})=-2(2^{d-1}-1)^2.
\]
Consequently, for $n=4d$ and every $\gamma\in[1/(4d^2),1/(2d^2))$,
$2K_{d,d}$ does not maximise $i_\gamma$, even among bipartite $d$-regular
graphs on $n$ vertices.
\end{thmBp}

So the answer is negative even within the bipartite class for which Kahn
proved the case $\gamma=0$. The graph $H_2$ is the cycle $C_8$.

At the exponential scale, on the other hand, the answer is positive for every
$\gamma$. For $q>0$ put
\[
 P_d(q)=\sum_{a,b=0}^d\binom da\binom db q^{ab},
\]
which is the sum of $q^{e(A)}$ over all subsets $A$ of $V(K_{d,d})$.

\begin{propC}
Let $d\ge2$ and $\gamma\ge0$, and put
\[
 \rho(d,\gamma)=\inf_{0<q\le1}\Bigl(\frac1{2d}\log P_d(q)-\gamma d\log q\Bigr).
\]
Every $d$-regular graph $G$ on $n$ vertices satisfies
$\frac1n\log i_\gamma(G)\le\rho(d,\gamma)$, and
$\frac1n\log i_\gamma(mK_{d,d})\to\rho(d,\gamma)$ as $n\to\infty$ through
multiples of $2d$, where $m=n/(2d)$.
\end{propC}

The upper bound is Carenini's Proposition~2.1 with the fugacity equal to~$1$,
optimised over~$q$. Proposition~C shows that the maximum of $i_\gamma$ over
$d$-regular graphs on $n$ vertices is $i_\gamma(mK_{d,d})\,e^{o(n)}$ for fixed
$d$ and $\gamma$; all the failures in this paper are subexponential. They are
not confined to small $n$, however.

\begin{thmD}
Let $d\ge2$ and $1/8<\gamma<1/2$. For $n$ divisible by $2d$ write
$n/(2d)=2k+r$ with $r\in\{0,1\}$, and let $G_n=kH_d\cup rK_{d,d}$, a bipartite
$d$-regular graph on $n$ vertices. Then $i_\gamma(G_n)>i_\gamma(mK_{d,d})$,
where $m=n/(2d)$, for all sufficiently large $n$ divisible by $2d$.
\end{thmD}

The proof compares the upper tails of the two counts with Cramér's theorem in
its qualitative form, so it gives no explicit threshold $n_0$ beyond which the
inequality holds.

\paragraph{Scope.}
Our counterexamples are exact-count effects. In Theorems~B and~B$'$ the
threshold is a single edge, $\lfloor\gamma dn\rfloor=1$. In Theorem~A,
$\gamma^*$ exceeds $1/8$ unless $(n,d)=(6,3)$, and $\gamma^*\to1/2$ as
$n\to\infty$. In Theorem~D, $\gamma>1/8$ is fixed; both counts are then
$(1-o(1))2^n$, and their relative difference tends to~$0$
(Remark~\ref{rem:size}). The value $1/8$ is where the mean crosses the
threshold: a uniformly random subset of a $d$-regular graph on $n$ vertices
spans $dn/8$ edges on average. Question~1.3 remains open for fixed
$\gamma\in(0,1/8]$ as $n\to\infty$. All our counterexamples with
$\gamma\le1/8$ have $\gamma d<1$. In particular, we know no counterexample
with $\gamma\le1/8$ and $\gamma d\ge2$, the range of small $\gamma$ in which
Carenini's Proposition~1.2 shows that the term $\gamma\log(e/\gamma)\,n$ of
her Theorem~1.1 is needed. Our counterexamples also do not refute the weaker
statement $i_\gamma(G)\le(1+o(1))\,i_\gamma(mK_{d,d})$.
Section~\ref{sec:readings} lists several readings of the question with our
status for each, and Section~\ref{sec:open} states open questions.

\paragraph{Prior work.}
Question~1.3 appeared in version~1 of~\cite{car}, submitted to the arXiv on
22 September 2026. We found no earlier answer to it, and no discussion of it,
in the following sources, which we searched on 26 September 2026 between
05:12 and 05:23 UTC: the arXiv (the listing of~\cite{car}, which had only
version~1, and searches of the arXiv API for later papers on the topic or by
the same author); Semantic Scholar, which listed no citations of~\cite{car};
the problem database MathDB (\url{https://mathdb.com}); GitHub searches of
code, commits, issues, pull requests and repositories; the public
repositories of AI-assisted work on open problems that we check, among them
the \texttt{trureturing} repository~\cite{trurepo}; and web search. Some of
these searches were repeated on 26 September 2026 between 06:39 and 06:42 UTC,
with the same result: we searched the issues and pull requests of the
\texttt{trureturing} repository for the versioned identifier
\texttt{2609.28527v1} (as a positive control, the same search for the
versioned identifier of another recent arXiv paper found the expected issue),
searched GitHub issues and commits, and searched MathDB by title. We did not
search Google Scholar, zbMATH or MathOverflow. This reports the coverage of
our searches, not a guarantee of priority. Theorem~1.1 and Proposition~2.1
of~\cite{car} are due to Carenini. The upper bound in Proposition~C is a
direct consequence of the latter, which rests on the theorem of Sah, Sawhney,
Stoner and Zhao~\cite[Corollary~1.15]{sssz}.

\paragraph{Organisation.}
Section~\ref{sec:readings} discusses readings of the question.
Section~\ref{sec:prelim} collects notation and a complement identity.
Theorem~A is proved in Section~\ref{sec:A}, Theorems~B and~B$'$ in
Section~\ref{sec:B}, Proposition~C in Section~\ref{sec:C} and Theorem~D in
Section~\ref{sec:D}. Section~\ref{sec:comp} reports computations,
Section~\ref{sec:lean} the formal verification, and Section~\ref{sec:open}
open questions.

\section{Readings of the question}\label{sec:readings}
Question~1.3 is a single universal statement, but it can be weakened in
several natural ways. Table~\ref{tab:readings} lists the readings we
considered; R1 is the literal one. In R4--R6, $d$ and $\gamma$ are fixed and
$n\to\infty$ through multiples of $2d$.

\begin{table}[ht]\centering\small
\begin{tabular}{@{}l >{\raggedright\arraybackslash}p{0.41\textwidth}
 >{\raggedright\arraybackslash}p{0.47\textwidth}@{}}\toprule
 & Reading & Status\\\midrule
R1 & Literal: for all admissible $(n,d)$ and all $\gamma\ge0$,
 $i_\gamma(G)\le i_\gamma(mK_{d,d})$ for every $d$-regular graph $G$ on $n$
 vertices.
 & False. Smallest counterexample (smallest $n$, then smallest $\gamma$):
 $(n,d,\gamma)=(6,3,1/18)$ (Theorems~A and~B); formally
 verified.\\\addlinespace
R2 & For a given admissible pair $(n,d)$, the inequality holds for all
 $\gamma\ge0$.
 & False for every admissible pair with $d\ge2$ and $(n,d)\ne(4,2)$
 (Theorem~A); true in the trivial cases $d=1$ and $(n,d)=(4,2)$.\\\addlinespace
R3 & Small $\gamma$: the inequality holds for all admissible $(n,d)$ and all
 $0\le\gamma<1/8$.
 & False: for $n=2d$ and every $d\ge3$ (Theorem~B); for $n=4d$ and every
 $d\ge2$, with a connected bipartite competitor (Theorem~B$'$); for some
 $n=2dm$ with $m\ge2$ (Corollary~\ref{cor:copies}); at $\gamma=1/10$ for many
 $n\le400$ (Remark~\ref{rem:tenth}).\\\addlinespace
R4 & Fixed $\gamma$: the inequality holds for all large $n$.
 & False for every $d\ge2$ and every $\gamma\in(1/8,1/2)$ (Theorem~D); true
 for $\gamma=0$ and for $\gamma\ge1/2$.\\\addlinespace
R5 & R4 for fixed $\gamma\in(0,1/8]$.
 & Open. True at the exponential scale (Proposition~C). The competitors of
 Theorem~D lose for large $n$ when $\gamma<1/8$
 (Remark~\ref{rem:lower}).\\\addlinespace
R6 & Ratio form: $\max_Gi_\gamma(G)\le(1+o(1))\,i_\gamma(mK_{d,d})$.
 & Not refuted. True for $\gamma=0$; for $\gamma=1/8$, where
 $i_{1/8}(G)=(1/2+o(1))2^n$ uniformly in $G$; and for $\gamma>1/8$, where
 $i_\gamma(mK_{d,d})=(1-o(1))2^n$ (Remark~\ref{rem:size}). Open for
 $0<\gamma<1/8$.\\\bottomrule
\end{tabular}
\caption{Readings of Question~1.3 of~\cite{car}; here $m=n/(2d)$.}
\label{tab:readings}
\end{table}

\FloatBarrier
\section{Preliminaries}\label{sec:prelim}
All graphs are finite and simple. For a graph $G$ let $N_s(G)$ be the number
of sets $A\subseteq V(G)$ with $e_G(A)=s$, and let
\[
 P_G(z)=\sum_{A\subseteq V(G)}z^{e_G(A)}=\sum_{s\ge0}N_s(G)\,z^s,
\]
so that $\Nle t(G)=\sum_{s\le t}[z^s]P_G(z)$ and $i_0(G)=N_0(G)$, where
$[z^s]F(z)$ denotes the coefficient of $z^s$ in a polynomial $F$. In the
notation of~\cite{car}, $P_G(q)=Z_G(1,q)$. For a disjoint union,
$P_{G\cup F}=P_GP_F$, since $e_{G\cup F}(A)=e_G(A\cap V(G))+e_F(A\cap V(F))$.
Choosing $a$ vertices on one side of $K_{d,d}$ and $b$ on the other spans
$ab$ edges, so $P_{K_{d,d}}=P_d$ and $P_{mK_{d,d}}=P_d^{\,m}$. For a
$d$-regular graph $G$ on $n$ vertices we write $E=e(G)=dn/2$. Every integer
$t\ge0$ arises as $\lfloor\gamma dn\rfloor$, namely for
$\gamma\in[t/(dn),(t+1)/(dn))$, so Question~1.3 is equivalent to the
inequalities $\Nle t(G)\le\Nle t(mK_{d,d})$ for all integers $t\ge0$.

\begin{lemma}[complement identity]\label{lem:complement}
Let $G$ be a $d$-regular graph with $n$ vertices and $E$ edges, and for
$S\subseteq V(G)$ let $\partial(S)$ be the number of edges meeting $S$. Then
$\partial(S)=d|S|-e_G(S)$, and for every integer $t$,
\[
 \Nle t(G)=2^n-\bigl|\{S\subseteq V(G):\ \partial(S)\le E-t-1\}\bigr|.
\]
\end{lemma}
\begin{proof}
Summing the degrees of the vertices of $S$ counts each edge inside $S$ twice
and each edge between $S$ and its complement once, so
$d|S|=e_G(S)+\partial(S)$. The edges inside $V(G)\setminus S$ are exactly the
edges not meeting $S$, so $e_G(V(G)\setminus S)=E-\partial(S)$. Hence
$e_G(A)\le t$ if and only if $\partial(V(G)\setminus A)\ge E-t$, and
$A\mapsto V(G)\setminus A$ is a bijection.
\end{proof}

\section{Proof of Theorem A}\label{sec:A}
\begin{proposition}\label{prop:formula}
Let $d\ge2$, let $G$ be a $d$-regular graph on $n$ vertices, and let
$t^*=dn/2-3d+1$. Then
\[
 \Nle{t^*}(G)=2^n-1-n-\binom n2-n\binom d2+2T(G)-Q(G),
\]
where $Q(G)$ is the number of $4$-sets spanning at least $d+2$ edges plus the
number of $5$-sets spanning at least $2d+2$ edges.
\end{proposition}
\begin{proof}
We have $E-t^*-1=3d-2$. By Lemma~\ref{lem:complement}, $2^n-\Nle{t^*}(G)$ is
the number of sets $S$ with $\partial(S)=d|S|-e_G(S)\le3d-2$. We count these
sets by their size $k$.
\begin{itemize}
\item $k\le2$: every set qualifies, since $\partial(S)\le kd\le2d\le3d-2$.
 This gives $1+n+\binom n2$ sets.
\item $k=3$: the condition is $e_G(S)\ge2$. Call a pair of edges with a
 common end a \emph{cherry}; there are $n\binom d2$ cherries. A $3$-set with
 exactly two edges contains exactly one cherry, and a triangle contains three.
 Hence there are $n\binom d2-3T(G)$ sets of size $3$ with exactly two edges,
 and $n\binom d2-2T(G)$ with at least two.
\item $k=4$ and $k=5$: the conditions are $e_G(S)\ge d+2$ and
 $e_G(S)\ge2d+2$, so these sets are counted by $Q(G)$.
\item $k\ge6$: we would need $e_G(S)\ge(k-3)d+2$, while $2e_G(S)\le kd$ as
 every vertex has degree $d$. Together these give $(k-6)d\le-4$, which is
 impossible.
\end{itemize}
Adding up proves the formula.
\end{proof}

\begin{lemma}\label{lem:Q}
Let $G$ be a $d$-regular graph. If $d\ge5$ then $Q(G)=0$; if $d=4$ then
$Q(G)\le\frac35T(G)$; if $d=3$ then $Q(G)\le\frac12T(G)$; and if $d=2$ then
$Q(G)$ is the number $c_4(G)$ of components of $G$ isomorphic to $C_4$. In
particular, if $d\ge3$ then $2T(G)-Q(G)\ge\frac75T(G)$.
\end{lemma}
\begin{proof}
If $d\ge5$, a $4$-set spans at most $6<d+2$ edges and a $5$-set spans at most
$10<2d+2$ edges.

Let $d=4$. Then $Q(G)$ is the number of copies of $K_4$ plus the number of
copies of $K_5$. A triangle $xyz$ lies in at most two copies of $K_4$, because
the fourth vertex is a neighbour of $x$ other than $y$ and $z$, and $x$ has
only two such neighbours. As each $K_4$ contains four triangles, the number
of copies of $K_4$ is at most $2T(G)/4$. Each vertex of a $K_4$ has exactly one
neighbour outside it, so a $K_4$ lies in at most one $K_5$. As each $K_5$
contains five copies of $K_4$, the number of copies of $K_5$ is at most a
fifth of the number of copies of $K_4$. Hence
$Q(G)\le\frac65\cdot\frac{T(G)}2=\frac35T(G)$.

Let $d=3$. A $5$-set spans at most $\lfloor15/2\rfloor=7<8$ edges, so $Q(G)$
counts the $4$-sets with at least five edges. Each of them contains at least
two triangles. A triangle $xyz$ lies in at most one of them: its fourth
vertex must be adjacent to at least two of $x,y,z$; if $w$ and $w'$ both have
this property, then some vertex of $\{x,y,z\}$ is adjacent to both, and since
that vertex has only one neighbour outside $\{x,y,z\}$, $w=w'$. Counting pairs
consisting of a triangle and a set counted by $Q(G)$ that contains it gives
$2Q(G)\le T(G)$.

Let $d=2$. A $4$-set with four edges induces a $2$-regular graph on four
vertices, that is, a component $C_4$; and a $5$-set spans at most five edges.

The last assertion follows because $\frac35\ge\frac12$.
\end{proof}

\begin{proof}[Proof of Theorem~A]
First, $t^*\ge1$: if $d\ge3$, then $n\ge2d$ gives $t^*\ge d^2-3d+1\ge1$; if
$d=2$, then $n\ge8$ and $t^*=n-5\ge3$. As $\gamma^*dn=t^*$, we have
$i_{\gamma^*}=\Nle{t^*}$. By Proposition~\ref{prop:formula}, for $d$-regular
graphs $G$ and $G'$ on $n$ vertices,
\[
 \Nle{t^*}(G)-\Nle{t^*}(G')=\bigl(2T(G)-Q(G)\bigr)-\bigl(2T(G')-Q(G')\bigr).
\]
The graph $mK_{d,d}$ is bipartite, so $T(mK_{d,d})=0$. If $d\ge3$, then also
$Q(mK_{d,d})=0$ by Lemma~\ref{lem:Q}, and (a) follows from
Lemma~\ref{lem:Q}: the difference $2T(G)-Q(G)$ is at least $\frac75T(G)$, and
it vanishes when $T(G)=0$. If $d=2$, Lemma~\ref{lem:Q} gives
$\Nle{t^*}(G)-\Nle{t^*}(mC_4)=2T(G)+n/4-c_4(G)$, which is positive unless
$T(G)=0$ and $c_4(G)=n/4$, that is, unless $G=mC_4$.
\end{proof}

\begin{remark}\label{rem:n2d}
For $n=2d$, every $d$-regular graph $G\ne K_{d,d}$ contains a triangle.
Indeed, if $G$ is triangle-free and $uv$ is an edge, the neighbourhoods $N(u)$
and $N(v)$ are disjoint independent sets of size $d$, so they partition
$V(G)$; every vertex of $N(u)$ then has all its $d$ neighbours in $N(v)$, and
$G=K_{d,d}$. (This is also the equality case of Mantel's
theorem~\cite{mantel}, since $G$ has $d^2$ edges.) Hence, for every $d\ge3$,
every $d$-regular graph on $2d$ vertices other than $K_{d,d}$ beats $K_{d,d}$
at $t^*=d^2-3d+1$. As a second
example, for $(n,d)=(12,3)$ we have $t^*=10$ and $\gamma^*=5/18$, and the
graph $3K_4$ has $T=12$ and $Q=3$; so
$\Nle{10}(3K_4)-\Nle{10}(2K_{3,3})=21$, the two counts being $4002$ and
$3981$.
\end{remark}

\section{\texorpdfstring{Small thresholds: proofs of Theorems B and B$'$}{Small thresholds: proofs of Theorems B and B'}}\label{sec:B}

\subsection{The switched complete bipartite graph}
\begin{proof}[Proof of Theorem~B]
Write $K=K_{d,d}$. Each of $x_1,x_2,y_1,y_2$ loses one edge and gains one, so
$S_d$ is $d$-regular. For $A\subseteq X\cup Y$ put $A_X=A\cap X$,
$A_Y=A\cap Y$, $a=|A_X|$ and $b=|A_Y|$, so that $e_K(A)=ab$.

\emph{Inclusion.} If $e_K(A)\le1$, then $a=0$, or $b=0$, or $a=b=1$. If
$b=0$, then $A\subseteq X$, and the only edge of $S_d$ inside $X$ is
$x_1x_2$; the case $a=0$ is symmetric; and a set of two vertices spans at most
one edge. So $e_{S_d}(A)\le1$. The inclusion is strict: $A_0=\{x_1,y_1,y_3\}$
has $e_K(A_0)=2$, but $e_{S_d}(A_0)=1$, since $x_1y_1$ was deleted and
$y_1y_3$ is not an edge. (Here $d\ge3$ is used.)

\emph{Exact count.} By the inclusion, $\Nle1(S_d)-\Nle1(K)$ is the number of
sets $A$ with $e_K(A)\ge2$ and $e_{S_d}(A)\le1$. Writing $[\cdot]$ for the
indicator of a statement,
\[
 e_{S_d}(A)=ab-[x_1\in A_X,\,y_1\in A_Y]-[x_2\in A_X,\,y_2\in A_Y]
 +[x_1,x_2\in A_X]+[y_1,y_2\in A_Y].
\]
We need $ab\ge2$ and $e_{S_d}(A)\le1$. At most two terms are subtracted, so
$ab\le3$.
\begin{itemize}
\item $(a,b)=(1,2)$: we need one subtraction and no addition, that is,
 $A_X=\{x_1\}$ and $y_1\in A_Y$, or $A_X=\{x_2\}$ and $y_2\in A_Y$, and in
 both cases $A_Y\ne\{y_1,y_2\}$. The second vertex of $A_Y$ is then one of
 $y_3,\dots,y_d$, so there are $2(d-2)$ such sets.
\item $(a,b)=(2,1)$: by symmetry, there are $2(d-2)$ such sets.
\item $ab=3$: then $a=1$ or $b=1$, at most one term is subtracted, and
 $e_{S_d}(A)\ge2$.
\end{itemize}
Hence $\Nle1(S_d)-\Nle1(K)=4d-8$. Finally, $\Nle1(K)=2^{d+1}-1+d^2$: the
independent sets of $K$ are the $2^{d+1}-1$ subsets of $X$ or of $Y$, and
the sets spanning exactly one edge are the $d^2$ pairs $\{x,y\}$ with $x\in X$
and $y\in Y$. For $n=2d$ and $\gamma\in[1/(2d^2),1/d^2)$ we have
$\lfloor\gamma dn\rfloor=1$, which gives the last assertion.
\end{proof}

\begin{corollary}\label{cor:copies}
Let $d\ge3$ and $m\ge2$, and put $I=2^{d+1}-1$. Then
\[
 \Nle1\bigl(S_d\cup(m-1)K_{d,d}\bigr)-\Nle1(mK_{d,d})
 =I^{m-2}\bigl[(4d-8)I-(m-1)d^2(2^{d-1}-2)\bigr].
\]
Hence, for $n=2dm$ and $\gamma\in[1/(dn),2/(dn))$, the graph
$S_d\cup(m-1)K_{d,d}$ beats $mK_{d,d}$ if and only if
$(4d-8)I>(m-1)d^2(2^{d-1}-2)$. This holds exactly for $d=3$ with
$2\le m\le4$, for $d\in\{4,5\}$ with $2\le m\le3$, and for $6\le d\le13$ with
$m=2$. For $d\ge14$ it fails for every $m\ge2$.
\end{corollary}
\begin{proof}
For disjoint graphs $G$ and $F$, a set spans at most one edge of $G\cup F$
exactly when it spans at most one edge of $G$ and none of $F$, or none of $G$
and exactly one of $F$. Hence
\[
 \Nle1(G\cup F)=\Nle1(G)\,i_0(F)+i_0(G)\bigl(\Nle1(F)-i_0(F)\bigr).
\]
Take $F=(m-1)K_{d,d}$, so that $i_0(F)=I^{m-1}$ and
$\Nle1(F)-i_0(F)=(m-1)d^2I^{m-2}$: one block spans one edge and the others
none. Subtracting the identities for $G=S_d$ and $G=K_{d,d}$ and using
Theorem~B, it remains to show $i_0(K_{d,d})-i_0(S_d)=2^{d-1}-2$, that is,
$i_0(S_d)=3\cdot2^{d-1}+1$. An independent set of $S_d$ inside $X$ is a subset
of $X$ not containing both $x_1$ and $x_2$, and there are $3\cdot2^{d-2}$ of
these; the same holds for $Y$. An independent set meeting both sides has
$a,b\ge1$; if $a=b=1$ it must be $\{x_1,y_1\}$ or $\{x_2,y_2\}$, and if
$ab\ge2$ then, by the formula for $e_{S_d}(A)$ above, $e_{S_d}(A)\ge ab-1\ge1$
when $\min(a,b)=1$, and $e_{S_d}(A)\ge ab-2\ge2$ otherwise. So
$i_0(S_d)=2\cdot3\cdot2^{d-2}-1+2=3\cdot2^{d-1}+1$.

For the list, note that the right-hand side of the criterion increases with
$m$, so for each $d$ the values of $m$ for which it holds form an interval
starting at $m=2$; for $3\le d\le13$ the list follows by direct evaluation.
For $d\ge14$, the right-hand side minus the left-hand side at $m=2$ equals
$(d^2-16d+32)2^{d-1}-2d^2+4d-8$, which is positive because
$d^2-16d+32=(d-8)^2-32\ge4$ and $2^{d+1}>2d^2+8$.
\end{proof}

\subsection{A two-block switch}
Let $K=K_{d,d}$ and, for $\alpha,\beta\in\{0,1\}$,
\[
 W_{\alpha\beta}(z)=\sum_{i,j=0}^{d-1}\binom{d-1}i\binom{d-1}j
 z^{ij+\alpha j+\beta i}.
\]
Exchanging $i$ and $j$ shows $W_{10}=W_{01}$. Put
$D(z)=W_{00}(z)W_{11}(z)-W_{10}(z)^2$.

\begin{lemma}\label{lem:switch}
Let $d\ge2$ and let $H_d$ be as in Theorem~B$'$. Then
\[
 P_{2K}(z)-P_{H_d}(z)=2(z-1)D(z).
\]
Consequently, for every graph $F$, possibly with no vertices, and every
integer $t\ge0$,
\[
 \Nle t(H_d\cup F)-\Nle t(2K\cup F)=2\,[z^t]\bigl(D(z)P_F(z)\bigr).
\]
\end{lemma}
\begin{proof}
The graphs $2K$ and $H_d$ have the same vertex set. For $A$ in it and $k=1,2$
let $\alpha_k=[u_k\in A]$, $\beta_k=[w_k\in A]$,
$i_k=|A\cap(X_k\setminus\{u_k\})|$ and $j_k=|A\cap(Y_k\setminus\{w_k\})|$.
The edges of the $k$-th copy of $K$ other than $u_kw_k$ that lie inside $A$
number $i_kj_k+\alpha_kj_k+\beta_ki_k$, and both graphs contain all of them.
In addition, $2K$ contains the edges $u_1w_1$ and $u_2w_2$, which contribute
$\alpha_1\beta_1+\alpha_2\beta_2$, while $H_d$ contains $u_1w_2$ and
$u_2w_1$, which contribute $\alpha_1\beta_2+\alpha_2\beta_1$. Summing over
$A$ with $(\alpha_1,\beta_1,\alpha_2,\beta_2)$ fixed,
\[
 P_{2K}=\sum W_{\alpha_1\beta_1}W_{\alpha_2\beta_2}\,
 z^{\alpha_1\beta_1+\alpha_2\beta_2},\qquad
 P_{H_d}=\sum W_{\alpha_1\beta_1}W_{\alpha_2\beta_2}\,
 z^{\alpha_1\beta_2+\alpha_2\beta_1},
\]
the sums being over $\{0,1\}^4$. The two exponents differ by
$(\alpha_1-\alpha_2)(\beta_1-\beta_2)$, which vanishes unless
$\alpha_1\ne\alpha_2$ and $\beta_1\ne\beta_2$. The patterns
$(\alpha_1,\beta_1,\alpha_2,\beta_2)=(1,1,0,0)$ and $(0,0,1,1)$ each
contribute $(z-1)W_{00}W_{11}$ to $P_{2K}-P_{H_d}$, and the patterns
$(1,0,0,1)$ and $(0,1,1,0)$ each contribute $(1-z)W_{10}W_{01}$. As
$W_{10}=W_{01}$, this proves the identity. For the second assertion, put
$R=DP_F$. Since $P_{G\cup F}=P_GP_F$,
\begin{align*}
 \Nle t(H_d\cup F)-\Nle t(2K\cup F)
 &=-2\sum_{s\le t}[z^s]\bigl((z-1)R\bigr)\\
 &=-2\sum_{s\le t}\bigl([z^{s-1}]R-[z^s]R\bigr)=2\,[z^t]R.\qedhere
\end{align*}
\end{proof}

\begin{lemma}\label{lem:sign}
Let $d\ge2$ and $q>0$. Then $D(q)>0$ if $q>1$, and $D(q)<0$ if $q<1$.
\end{lemma}
\begin{proof}
For $i,j\in\{0,\dots,d-1\}$ put
$\mu(i,j)=\binom{d-1}i\binom{d-1}jq^{ij}$, and consider
\[
 \Delta=\sum_{i,i',j,j'}\bigl[\mu(i,j)\mu(i',j')-\mu(i,j')\mu(i',j)\bigr]
 \bigl(q^i-q^{i'}\bigr)\bigl(q^j-q^{j'}\bigr),
\]
where all four indices range over $\{0,\dots,d-1\}$. We have
$W_{00}(q)=\sum\mu(i,j)$, $W_{11}(q)=\sum\mu(i,j)q^{i+j}$ and
$W_{10}(q)=\sum\mu(i,j)q^j=\sum\mu(i,j)q^i$. Expanding
$(q^i-q^{i'})(q^j-q^{j'})$ into four terms, we find
\begin{align*}
 \sum_{i,i',j,j'}\mu(i,j)\mu(i',j')\bigl(q^i-q^{i'}\bigr)\bigl(q^j-q^{j'}\bigr)
 &=2W_{00}W_{11}-2W_{10}^2,\\
 \sum_{i,i',j,j'}\mu(i,j')\mu(i',j)\bigl(q^i-q^{i'}\bigr)\bigl(q^j-q^{j'}\bigr)
 &=2W_{10}^2-2W_{00}W_{11},
\end{align*}
all at the point $q$. Subtracting, $\Delta=4D(q)$.

Now let $q\ne1$. The bracket equals
$\binom{d-1}i\binom{d-1}{i'}\binom{d-1}j\binom{d-1}{j'}\,
q^{ij'+i'j}\bigl(q^{(i-i')(j-j')}-1\bigr)$, which has the sign of
$(i-i')(j-j')\log q$. Since $q^x$ is strictly monotone in $x$, the product
$(q^i-q^{i'})(q^j-q^{j'})$ has the sign of $(i-i')(j-j')$. So every term of
$\Delta$ is zero or has the sign of $\log q$, and the term with $i=j=0$ and
$i'=j'=1$, which exists because $d-1\ge1$, is not zero.
\end{proof}

\begin{proof}[Proof of Theorem~B$'$]
Each of $u_1,u_2,w_1,w_2$ loses one edge and gains one, so $H_d$ is
$d$-regular. It is bipartite with sides $X_1\cup X_2$ and $Y_1\cup Y_2$,
since the new edges join $X_1$ to $Y_2$ and $X_2$ to $Y_1$. A copy of
$K_{d,d}$ minus one edge is connected when $d\ge2$, and the edge $u_1w_2$ joins
the two copies, so $H_d$ is connected.

By Lemma~\ref{lem:switch} with $F$ the graph with no vertices ($P_F=1$),
$\Nle t(H_d)-\Nle t(2K_{d,d})=2[z^t]D$ for every $t\ge0$. The exponent in
$W_{00}$ is $ij$, so $[z^0]W_{00}=2^d-1$ (the pairs with $i=0$ or $j=0$) and
$[z^1]W_{00}=(d-1)^2$. The exponent in $W_{11}$ is $(i+1)(j+1)-1$, so
$[z^0]W_{11}=1$ and $[z^1]W_{11}=2(d-1)$. The exponent in $W_{10}$ is
$j(i+1)$, so $[z^0]W_{10}=2^{d-1}$ and $[z^1]W_{10}=d-1$. Hence
\begin{align*}
 [z^0]D&=2^d-1-4^{d-1}=-(2^{d-1}-1)^2,\\
 [z^1]D&=2(d-1)(2^d-1)+(d-1)^2-2\cdot2^{d-1}(d-1)=(d-1)(2^d+d-3).
\end{align*}
Since $i_0=N_{\le0}$, this proves the two formulas. For $n=4d$ and
$\gamma\in[1/(4d^2),1/(2d^2))$ we have $\lfloor\gamma dn\rfloor=1$.
\end{proof}

\section{The exponential scale: proof of Proposition C}\label{sec:C}
Let $\xi$ be the number of edges of $K_{d,d}$ spanned by a uniformly random
subset of $V(K_{d,d})$. Then $\xi$ takes values in $\{0,\dots,d^2\}$,
$\Prb(\xi=0)>0$, $\Prb(\xi=d^2)=4^{-d}$, $\Ex\xi=d^2/4$ because each edge is
spanned with probability $1/4$, and $\Ex q^\xi=4^{-d}P_d(q)$ for $q>0$. For
$m\ge1$ let $\Sigma_m=\xi_1+\dots+\xi_m$, where $\xi_1,\dots,\xi_m$ are
independent copies of $\xi$. A uniformly random subset of $V(mK_{d,d})$ meets
the $m$ copies in
independent uniformly random subsets. Hence, for $n=2dm$ and $x=2\gamma d^2$,
so that $\gamma dn=mx$,
\begin{equation}\label{eq:prob}
 i_\gamma(mK_{d,d})=2^n\,\Prb(\Sigma_m\le mx),\qquad
 2^n-i_\gamma(mK_{d,d})=2^n\,\Prb(\Sigma_m>mx).
\end{equation}
We use the lower bound in Cramér's theorem in the following form.

\begin{lemma}\label{lem:cramer}
Let $\xi$ be a random variable with values in $\{0,1,\dots,M\}$ such that
$\Prb(\xi=0)>0$ and $\Prb(\xi=M)>0$, let $\Sigma_m$ be the sum of $m$ independent
copies of $\xi$, and let $\Lambda(\lambda)=\log\Ex e^{\lambda\xi}$.
\begin{itemize}
\item[(a)] If $\Ex\xi<x<M$, then
 $\liminf_{m\to\infty}\frac1m\log\Prb(\Sigma_m>mx)\ge
 \inf_{\lambda\ge0}\bigl(\Lambda(\lambda)-\lambda x\bigr)$.
\item[(b)] If $0<x<\Ex\xi$, then
 $\liminf_{m\to\infty}\frac1m\log\Prb(\Sigma_m<mx)\ge
 \inf_{\lambda\le0}\bigl(\Lambda(\lambda)-\lambda x\bigr)$.
\end{itemize}
\end{lemma}
\begin{proof}
Let $\Lambda^*(y)=\sup_{\lambda\in\mathbb R}\bigl(\lambda y-\Lambda(\lambda)\bigr)$.
By Cramér's theorem~\cite[Theorem~2.2.3]{dz}, for every open set
$U\subseteq\mathbb R$,
\[
 \liminf_{m\to\infty}\frac1m\log\Prb(\Sigma_m/m\in U)\ge-\inf_{y\in U}\Lambda^*(y).
\]
The function $\Lambda^*$ is convex, as a supremum of affine
functions. It is finite on $[0,M]$: for all $\lambda$ we have
$\Lambda(\lambda)\ge\lambda M+\log\Prb(\xi=M)$ and
$\Lambda(\lambda)\ge\log\Prb(\xi=0)$, so for $0\le y\le M$ the quantity
$\lambda y-\Lambda(\lambda)$ is at most $-\log\Prb(\xi=M)$ when $\lambda\ge0$ and
at most $-\log\Prb(\xi=0)$ when $\lambda\le0$. Hence $\Lambda^*$ is continuous on
$(0,M)$. By Jensen's inequality $\Lambda(\lambda)\ge\lambda\Ex\xi$, so for
$y\ge\Ex\xi$ the terms with $\lambda<0$ are at most $0$, the value at
$\lambda=0$; thus $\Lambda^*(y)=\sup_{\lambda\ge0}(\lambda y-\Lambda(\lambda))$, which
is nondecreasing in $y\ge\Ex\xi$. In case~(a), take $U=(x,\infty)$; then
$\inf_{y>x}\Lambda^*(y)=\lim_{y\downarrow x}\Lambda^*(y)=\Lambda^*(x)=
-\inf_{\lambda\ge0}(\Lambda(\lambda)-\lambda x)$. Case~(b) is symmetric, with
$U=(-\infty,x)$ and $\Lambda^*$ nonincreasing on $(-\infty,\Ex\xi]$.
\end{proof}

\begin{proof}[Proof of Proposition~C]
\emph{Upper bound.} By~\cite[Proposition~2.1]{car} with $\lambda=1$, every
$d$-regular graph $G$ on $n$ vertices satisfies
$i_\gamma(G)\le e^{\gamma tn}P_d(e^{-t/d})^{n/(2d)}$ for every $t\ge0$. With
$q=e^{-t/d}\in(0,1]$ this reads $i_\gamma(G)\le q^{-\gamma dn}P_d(q)^{n/(2d)}$.
Taking logarithms, dividing by $n$ and taking the infimum over $q$ gives
$\frac1n\log i_\gamma(G)\le\rho(d,\gamma)$. (The proof in~\cite{car} combines
the bound $P_G(q)\le P_d(q)^{n/(2d)}$ of~\cite[Corollary~1.15]{sssz} with the
inequality $[e_G(A)\le\gamma dn]\le q^{e_G(A)-\gamma dn}$.)

\emph{Lower bound.} Let $n=2dm$, and let $x=2\gamma d^2$ and $\Sigma_m$ be as
in~\eqref{eq:prob}. By the upper bound, it suffices to show that the lower
limit of $\frac1n\log i_\gamma(mK_{d,d})$ as $m\to\infty$ is at least
$\rho(d,\gamma)$.
\begin{itemize}
\item $\gamma=0$: here $i_0(mK_{d,d})=i_0(K_{d,d})^m$ exactly, so
 $\frac1n\log i_0(mK_{d,d})=\frac1{2d}\log i_0(K_{d,d})=
 \lim_{q\to0}\frac1{2d}\log P_d(q)\ge\rho(d,0)$.
\item $0<\gamma<1/8$: then $0<x<d^2/4=\Ex\xi$. By~\eqref{eq:prob} and
 Lemma~\ref{lem:cramer}(b), with $q=e^\lambda$,
 \begin{align*}
 \liminf_{m\to\infty}\frac1m\log i_\gamma(mK_{d,d})
 &\ge2d\log2+\inf_{0<q\le1}\bigl(\log(4^{-d}P_d(q))-x\log q\bigr)\\
 &=\inf_{0<q\le1}\bigl(\log P_d(q)-x\log q\bigr).
 \end{align*}
 Dividing by $2d$ and noting $x/(2d)=\gamma d$ gives the claim.
\item $\gamma\ge1/8$: then $x\ge\Ex\xi$, and
 $\Prb(\Sigma_m\le mx)\ge\Prb(\Sigma_m\le m\Ex\xi)\to1/2$ by the central limit theorem,
 as $\xi$ is not constant. Hence $\frac1n\log i_\gamma(mK_{d,d})\to\log2$, and
 $\log2\ge\rho(d,\gamma)$ is the value at $q=1$.\qedhere
\end{itemize}
\end{proof}

\begin{remark}
In the notation of~\cite{car}, $P_d(e^{-t/d})=B_d(t)$, so
$\rho(d,\gamma)=\inf_{t\ge0}\bigl(\gamma t+\frac1{2d}\log B_d(t)\bigr)$.
Carenini's Theorem~1.1 bounds this infimum by choosing $t=2\log(e/\gamma)$.
\end{remark}

\section{Fixed thresholds above one eighth: proof of Theorem D}\label{sec:D}
Let $n=2dm$, $x=2\gamma d^2$ and $\xi$, $\Sigma_m$ as in Section~\ref{sec:C}; for
$1/8<\gamma<1/2$ we have $d^2/4<x<d^2$. For a graph $G$ on $n$ vertices let
$U(G)=2^n-i_\gamma(G)$, the number of sets spanning more than $\gamma dn$
edges. We must show $U(G_n)<U(mK_{d,d})$ for all large $m$. Write $H=H_d$.

\begin{proof}[Proof of Theorem~D]
\emph{Lower bound for $U(mK_{d,d})$.} Let $\varphi(q)=\log P_d(q)-x\log q$
for $q>0$. Since $P_d(q)=4^d\Ex q^\xi$, we have $\varphi'(1)=\Ex\xi-x<0$;
and $\varphi(q)\ge(d^2-x)\log q\to\infty$ as $q\to\infty$, because
$P_d(q)\ge q^{d^2}$. Hence $\varphi$ attains its minimum over $[1,\infty)$ at
some $q^*>1$. By~\eqref{eq:prob} and Lemma~\ref{lem:cramer}(a), with
$q=e^\lambda$, for every $\varepsilon>0$ and all large $m$,
\[
 U(mK_{d,d})=2^n\Prb(\Sigma_m>mx)\ge2^n\exp\bigl(m(\varphi(q^*)-d\log4-\varepsilon)
 \bigr)=P_d(q^*)^m(q^*)^{-mx}e^{-\varepsilon m}.
\]

\emph{Upper bound for $U(G_n)$.} For $q\ge1$ and every set $A$ we have
$[e(A)>\gamma dn]\le q^{e(A)-\gamma dn}$. Summing over $A$ with $q=q^*$ gives
$U(G_n)\le(q^*)^{-mx}P_{G_n}(q^*)=(q^*)^{-mx}P_H(q^*)^kP_d(q^*)^r$.

\emph{Comparison.} As $m-r=2k$,
\[
 \frac{U(mK_{d,d})}{U(G_n)}\ge
 \Bigl(\frac{P_d(q^*)^2}{P_H(q^*)}\Bigr)^ke^{-\varepsilon m}.
\]
By Lemma~\ref{lem:switch} at $z=q^*$ and Lemma~\ref{lem:sign},
$P_d(q^*)^2-P_H(q^*)=2(q^*-1)D(q^*)>0$, so
$\beta=\log\bigl(P_d(q^*)^2/P_H(q^*)\bigr)>0$. Take $\varepsilon=\beta/4$.
Since $k\ge(m-1)/2$, the ratio is at least
$\exp\bigl(\beta(m-1)/2-\beta m/4\bigr)=\exp\bigl(\beta(m-2)/4\bigr)>1$ for all
large $m$. Hence $U(G_n)<U(mK_{d,d})$, that is,
$i_\gamma(G_n)>i_\gamma(mK_{d,d})$, for all large $n$ divisible by $2d$.
\end{proof}

\begin{remark}\label{rem:size}
For $\gamma>1/8$ the law of large numbers gives $\Prb(\Sigma_m\le mx)\to1$, so
$i_\gamma(mK_{d,d})=(1-o(1))2^n$, and the relative difference
$i_\gamma(G_n)/i_\gamma(mK_{d,d})-1\le U(mK_{d,d})/i_\gamma(mK_{d,d})$ tends to
$0$ exponentially fast. For example, for $d=2$, $\gamma=1/4$ and $n=400$ the
relative difference is about $1.3\cdot10^{-17}$. In particular, Theorem~D is
consistent with reading R6. Our proof gives no explicit $n_0$; one could be
extracted from a quantitative form of Cramér's theorem (for example,
exponential tilting combined with the Berry--Esseen theorem, or the theorem of
Bahadur and Rao~\cite{br}), which we have not done. The base
$P_d(q^*)^2/P_H(q^*)$ tends to $1$ as $\gamma\downarrow1/8$, since then
$q^*\to1$; for $d=2$ it is about $1.0006$ at $\gamma=1/4$ and $1.0001$ at
$\gamma=0.49$. So $n_0$ can be large, and we give no value.

Reading R6 also holds at $\gamma=1/8$. Let $A$ be a uniformly random subset of
$V(G)$, where $G$ is $d$-regular on $n$ vertices. Then $e_G(A)$ is the sum of
the indicators $[u,v\in A]$ over the $dn/2$ edges $uv$ of $G$; it has mean
$dn/8$ and variance $dn(2d+1)/32$, as each indicator has variance $3/16$ and
each of the $nd(d-1)$ ordered pairs of distinct edges with a common end has
covariance $1/16$. The indicators of two sets of edges that share no vertex are
independent, so the line graph of $G$, of maximum degree $2d-2$, is a
dependency graph for them, and a normal approximation
bound~\cite[Theorem~3.5 and Proposition~1.2]{ross} gives
$i_{1/8}(G)=2^n\Prb(e_G(A)\le dn/8)=(1/2+O(n^{-1/4}))\,2^n$, uniformly in $G$
for fixed $d$.
\end{remark}

\begin{remark}\label{rem:lower}
For $0<\gamma<1/8$ the same argument, applied to the lower tail, shows that
$G_n$ loses to $mK_{d,d}$ for all large $n$. Indeed, now $0<x<\Ex\xi$, so
$\varphi$ is minimised over $(0,1]$ at some $q^*<1$ (note $\varphi'(1)>0$ and
$\varphi(q)\to\infty$ as $q\to0$). By Lemma~\ref{lem:sign},
$D(q^*)<0$, so again $P_d(q^*)^2-P_H(q^*)=2(q^*-1)D(q^*)>0$.
Lemma~\ref{lem:cramer}(b) gives
$i_\gamma(mK_{d,d})\ge P_d(q^*)^m(q^*)^{-mx}e^{-\varepsilon m}$, the inequality
$[e(A)\le\gamma dn]\le q^{e(A)-\gamma dn}$ for $q\le1$ gives
$i_\gamma(G_n)\le(q^*)^{-mx}P_H(q^*)^kP_d(q^*)^r$, and the comparison is as
above. This is consistent with Proposition~C.
\end{remark}

\begin{remark}\label{rem:tenth}
The small-$\gamma$ failures are nevertheless not confined to tiny graphs. At
$\gamma=1/10$, the graph $G_n=kH_d\cup rK_{d,d}$ of Theorem~D beats
$mK_{d,d}$ for $92$ of the $119$ values $n\in\{8,12,\dots,480\}$ when $d=2$,
the largest being $n=400$, and for $56$ of the $119$ values
$n\in\{12,18,\dots,720\}$ when $d=3$, the largest being $n=360$
(Section~\ref{sec:comp}). By Remark~\ref{rem:lower} these are finite-size
effects.
\end{remark}

\section{Computations}\label{sec:comp}
All counts in this paper were computed in exact integer arithmetic. Our
programs compute the distribution $(N_s(G))_s$ in two independent ways: by a
recursion over bitmasks, $e_G(A\cup\{v\})=e_G(A)+|N(v)\cap A|$ for a vertex
$v$ larger than every element of $A$, and by a direct loop over all sets and
all edges. The two agree on eleven test graphs. Counts for disjoint unions are
obtained from $P_{G\cup F}=P_GP_F$. Two referee agents (see the disclosure at
the end of the paper) recomputed the main numbers, each with its own programs.
Our programs and their logs are available in version~1.3.0 of the
repository~\cite{repo} (directory \texttt{certificates/carenini/}).

\paragraph{Exhaustive search at $n=2d$.}
For $d\in\{3,4,5\}$ we generated all $d$-regular graphs on $2d$ vertices up to
isomorphism, starting from $K_{d,d}$ and applying $2$-switches (replacing two
edges $ab$, $cd$ by $ac$, $bd$ or by $ad$, $bc$ when these are not edges), with
exact isomorphism tests. Any two graphs with the same degree sequence are
connected by a sequence of $2$-switches~\cite{hakimi,fhm}, so this finds every
class. We found $2$, $6$ and $60$ classes. Every $d$-regular graph on $2d$
vertices is connected, since each component has at least $d+1$ vertices, and
these are the numbers of connected cubic graphs on $6$ vertices, quartic
graphs on $8$ vertices and quintic graphs on $10$ vertices (sequences A002851,
A006820 and A006821 of~\cite{oeis}). The first referee agent confirmed the cases
$d=3,4$ from all $70$ and $19355$ labelled graphs. Table~\ref{tab:exh} lists
the thresholds at which $K_{d,d}$ is not a maximiser.

\begin{table}[ht]\centering\small
\begin{tabular}{@{}cccccc@{}}\toprule
$d$ & $n$ & classes & $t$ & $\max_G\Nle t(G)$ & $\Nle t(K_{d,d})$\\\midrule
3 & 6 & 2 & 1 & 28 & 24\\\addlinespace
4 & 8 & 6 & 1 & 61 & 47\\
 & & & 3 & 137 & 127\\
 & & & 5 & 187 & 171\\\addlinespace
5 & 10 & 60 & 1 & 116 & 88\\
 & & & 2 & 203 & 188\\
 & & & 3 & 321 & 288\\
 & & & 5 & 524 & 448\\
 & & & 7 & 679 & 648\\
 & & & 8 & 783 & 748\\
 & & & 11 & 908 & 868\\\bottomrule
\end{tabular}
\caption{All thresholds $t$ at which $K_{d,d}$ does not maximise $\Nle t$
among $d$-regular graphs on $n=2d$ vertices, for $3\le d\le5$. The maximum is over
all such graphs; the failure holds for $\gamma\in[t/(2d^2),(t+1)/(2d^2))$.}
\label{tab:exh}
\end{table}

At every other threshold $t\le d^2$, $K_{d,d}$ is a maximiser. It is the unique
maximiser for $t\in\{0,2\}$ when $d=3$, for $t\in\{0,4,6\}$ when $d=4$, and
for $t\in\{0,4,6,9,10,12\}$ when $d=5$; at the remaining thresholds it ties
with other graphs. (By Lemma~\ref{lem:complement}, for $d\ge3$ all $d$-regular
graphs on $n$ vertices have the same value of $\Nle t$ when $t\ge dn/2-2d$.)
The failures at $t=1$ are those of Theorem~B, and the failures at
$t=d^2-3d+1$, that is, $t=1,5,11$, are those of Theorem~A (see also
Remark~\ref{rem:n2d}); the other failures are not covered by our theorems.

\paragraph{Other checks.}
The formula of Proposition~\ref{prop:formula} was checked on $100$ regular
graphs with $2\le d\le8$ and $n\le20$ (unions of copies of $K_{d,d}$,
$(K_d\cart K_2)\cup(m-1)K_{d,d}$, cycles, $C_3\cup C_{n-3}$, unions of copies
of $K_{d+1}$, and random regular graphs), with no mismatch; the first referee
agent checked $75$ further graphs, including graphs in which $5$-sets contribute
to~$Q$. Theorem~B was checked for $3\le d\le12$, and
Corollary~\ref{cor:copies} against brute force for
$(d,m)\in\{(3,2),(3,3),(4,2),(5,2)\}$. The identity of
Lemma~\ref{lem:switch} was checked symbolically for $2\le d\le5$, and the
difference in Theorem~B$'$ for $2\le d\le20$; the sign in
Lemma~\ref{lem:sign} was checked on a grid of values of $q$, as a sanity
check. For Theorem~D, we compared $kH_d$ with $2kK_{d,d}$ exactly for
$d\in\{2,3,4\}$ and all $n\le600$ divisible by $4d$. At $n=192$ the competitor
wins for $\gamma\in\{1/5,1/4,1/3,2/5,9/20\}$, while for $\gamma=3/20$, closer
to $1/8$, it still loses, as expected when $n_0$ is large. The outcome is not
monotone in $n$, so this does not show that $n_0\le192$: for $d=4$ and
$\gamma=9/20$ the competitor loses at $n=128$ and again at $n=208$, and for
$d=3$ and $\gamma=9/20$ it wins at $n=96$ and $n=192$ but loses at $n=144$.
For $\gamma=3/20$ it wins for every such $n$ from $360$ to $600$ when $d=2$,
from $336$ to $600$ when $d=3$, and from $320$ to $592$ when $d=4$, in line
with Theorem~D. The counts in
Remark~\ref{rem:tenth} are exact comparisons for all $n$ in the stated
ranges.

\section{Formal verification}\label{sec:lean}
The Lean file \texttt{lean/Research/Carenini.lean} in version~1.3.0 of the
repository~\cite{repo} formalises Question~1.3 literally and proves five
negative results: four instances and one infinite family ($d\ge3$). It uses
Lean~4
(version 4.33.1) and Mathlib~\cite{mathlib} (commit \texttt{0df444a3}); all
declarations are in the namespace \texttt{Carenini}. The basic definitions are
the following.
\begin{itemize}
\item \texttt{edgesIn G A} is the number of edges of \texttt{G} (elements of
 Mathlib's \texttt{G.edgeFinset}) with both ends in the finite set \texttt{A};
 this is $e_G(A)$. \texttt{iCount G t} is the number of sets \texttt{A} with
 \texttt{edgesIn G A} $\le t$, for $t\in\mathbb N$; this is $\Nle t(G)$.
\item \texttt{iGamma G d gamma} is the number of sets \texttt{A} with
 $\texttt{edgesIn G A}\le\gamma\cdot d\cdot|V|$, the comparison being made in
 $\mathbb R$ for a real number $\gamma$; this is $i_\gamma(G)$.
\item \texttt{KddUnion m d} is the graph on
 $\texttt{Fin m}\times(\texttt{Fin d}\oplus\texttt{Fin d})$ in which $(i,x)$
 and $(j,y)$ are adjacent when $i=j$ and $x,y$ are adjacent in Mathlib's
 \texttt{completeBipartiteGraph (Fin d) (Fin d)}; this is $mK_{d,d}$.
\end{itemize}
The question and the two main theorems read, in ASCII notation:
\begin{verbatim}
def CareniniQuestion (n d : Nat) (gamma : Real) : Prop :=
  forall (V : Type) [Fintype V] [DecidableEq V] (G : SimpleGraph V)
    [DecidableRel G.Adj], Fintype.card V = n -> G.IsRegularOfDegree d ->
      iGamma G d gamma <= iGamma (KddUnion (n / (2 * d)) d) d gamma

theorem not_careniniQuestion_six : Not (CareniniQuestion 6 3 (1 / 18))

theorem not_careniniQuestion_two_d {d : Nat} (hd : 3 <= d) :
    Not (CareniniQuestion (2 * d) d (1 / (2 * (d : Real) ^ 2)))
\end{verbatim}
The predicate quantifies over every finite vertex type of cardinality $n$ and
every $d$-regular simple graph on it (Mathlib's \texttt{IsRegularOfDegree});
the threshold $\gamma dn$ is a real number, as in the question; and the
reference graph consists of $n/(2d)$ disjoint copies of Mathlib's $K_{d,d}$.
The standing hypothesis $2d\mid n$ is not part of the predicate. Every
instance refuted below satisfies it, and at such an instance the predicate is
exactly the claim of Question~1.3. The lemma \texttt{iGamma\_eq\_iCount\_floor}
shows that $i_\gamma=\Nle{\lfloor\gamma dn\rfloor}$ for $\gamma\ge0$, and
\texttt{iGamma\_iso} shows that $i_\gamma$ is invariant under isomorphism.
Table~\ref{tab:lean} lists the five negative results: four instances and one
infinite family ($d\ge3$).

\begin{table}[ht]\centering\small
\begin{tabular}{@{}lll@{}}\toprule
Lean theorem & $(n,d,\gamma)$ & competitor and counts\\\midrule
\texttt{not\_careniniQuestion\_six} & $(6,3,1/18)$ &
 prism, $28>24$ (Theorems~A, B)\\
\texttt{not\_careniniQuestion\_two\_d} & $(2d,d,1/(2d^2))$, $d\ge3$ &
 $S_d$ (Theorem~B)\\
\texttt{not\_careniniQuestion\_eight\_two} & $(8,2,1/16)$ &
 $C_8=H_2$, $111>105$ (Theorem~B$'$)\\
\texttt{not\_careniniQuestion\_twelve\_bip} & $(12,3,1/36)$ &
 $H_3$, $527>495$ (Theorem~B$'$)\\
\texttt{not\_careniniQuestion\_twelve\_top} & $(12,3,5/18)$ &
 $3K_4$, $4002>3981$ (Theorem~A)\\\bottomrule
\end{tabular}
\caption{The negative results proved in Lean. In each case the reference
graph is $\frac n{2d}K_{d,d}$.}\label{tab:lean}
\end{table}

Supporting declarations make the statements concrete. For $n=6$,
\texttt{prism\_regular} and \texttt{K33\_regular} are proved by
\texttt{decide}, the counts $28$ and $24$ are evaluated by \texttt{decide},
and \texttt{prismIso}, \texttt{K33Iso} and \texttt{K33IsoKddUnion} show that the
two graphs are Mathlib's $K_3\cart K_2$ and $K_{3,3}$, the latter also as
\texttt{KddUnion 1 3}; \texttt{prism\_beats\_K33} is the explicit form of the
instance. For general $d$, \texttt{switchedKdd\_regular} shows that $S_d$ is
$d$-regular, \texttt{edgesIn\_switched\_le\_one} is the inclusion in
Theorem~B, \texttt{switched\_extra\_set} is the witness $\{x_1,y_1,y_3\}$
(indexed from $0$ in Lean), \texttt{iCount\_one\_lt} gives
$\Nle1(K_{d,d})<\Nle1(S_d)$, and \texttt{kddUnionOneIso} identifies $K_{d,d}$
with \texttt{KddUnion 1 d}. Finally, \texttt{H3\_regular},
\texttt{H3\_bipartite} and \texttt{H3\_connected} prove the properties of $H_3$
stated in Theorem~B$'$.

The following are not formalised: Theorem~A (only two instances of its
conclusion are formal, the prism at $(6,3,1/18)$ and $3K_4$ at $(12,3,5/18)$;
Proposition~\ref{prop:formula}, Lemma~\ref{lem:Q} and the minimiser statement
are not); the exact value $4d-8$ in Theorem~B (only the strict inequality is)
and Corollary~\ref{cor:copies}; Theorem~B$'$ (only
$\Nle1(H_d)>\Nle1(2K_{d,d})$ for $d=2,3$, via the counts $111>105$ and
$527>495$, the regularity of $H_2=C_8$ and of $H_3$, and the bipartiteness and
connectivity of $H_3$ are formal), and Lemmas~\ref{lem:switch}
and~\ref{lem:sign}; Proposition~C and
Theorem~D, which use Cramér's theorem and, for the upper bound in
Proposition~C, the theorem of Sah, Sawhney, Stoner and Zhao; and the
computations of Section~\ref{sec:comp}.

The file contains no \texttt{sorry} and does not use \texttt{native\_decide};
the finite counts are evaluated by the kernel, using \texttt{decide} and
\texttt{decide +kernel}. The file was rebuilt with \texttt{lake build} in a
clean directory, with exit code~$0$, and the compiled module was checked again
with \texttt{leanchecker}. The recorded axiom audits (\texttt{\#print axioms})
of the fourteen main declarations report only the standard axioms
\texttt{propext}, \texttt{Classical.choice} and \texttt{Quot.sound}, and no
axioms at all for \texttt{H3\_bipartite}. An audit file,
\texttt{lean/Research/CareniniAudit.lean}, prints the basic definitions and the
statements of Table~\ref{tab:lean} together with their axioms. The Lean file,
the audit file and the logs of the clean rebuild and of the \texttt{leanchecker}
run (files \texttt{logs/clean-replay-v1.3.0-*}) are in version~1.3.0 of the
repository~\cite{repo}.

\section{Open questions}\label{sec:open}
\begin{enumerate}
\item (Reading R5.) Let $d\ge2$ and $0<\gamma\le1/8$ be fixed. Is
 $i_\gamma(G)\le i_\gamma(mK_{d,d})$ for every $d$-regular graph $G$ on $n$
 vertices, for all sufficiently large $n$ divisible by $2d$? At
 $\gamma=1/10$ the inequality fails for many $n$ up to $400$
 (Remark~\ref{rem:tenth}), but the competitors used there lose for large $n$
 (Remark~\ref{rem:lower}). For $\gamma<1/8$ the infimum defining
 $\rho(d,\gamma)$ is attained at some $q<1$, in the antiferromagnetic range in
 which the theorem of Sah, Sawhney, Stoner and Zhao favours $K_{d,d}$. This
 suggests a positive answer, but we make no claim.
\item All our counterexamples with $\gamma\le1/8$ have $\gamma d<1$. Is
 there a counterexample, for some $n$, with $\gamma\le1/8$ and $\gamma d\ge2$,
 the range of small $\gamma$ in Carenini's Proposition~1.2?
\item (Reading R6.) For fixed $d$ and $0<\gamma<1/8$, is the maximum of
 $i_\gamma$ over $d$-regular graphs on $n$ vertices at most
 $(1+o(1))\,i_\gamma(mK_{d,d})$ as $n\to\infty$? (For $\gamma=1/8$ it is; see
 Remark~\ref{rem:size}.)
\item Which $d$-regular graphs maximise $i_\gamma$ when $1/8<\gamma<1/2$ and
 $n$ is large? By Proposition~\ref{prop:formula}, at $\gamma=\gamma^*$ the
 maximisers are the graphs that maximise $2T(G)-Q(G)$.
\end{enumerate}

\section*{Tool and computational resource disclosure}
This work was carried out on 26 September 2026 (UTC) in a workflow directed
by the author, using AI coding agents: Anthropic Claude Code, with the model
Claude Opus~5.5. The author set the research programme, namely resolving
open problems with AI agents and formal verification, and directed the
agents; the agents selected this question. Claude Code agents found
the counterexamples and the proofs, wrote the verification programs and the
Lean formalisation, carried out the literature searches described in the
introduction, and drafted this text. An independent Claude Code agent, without
access to the reasoning of the agents that found the proofs, refereed the
proofs and the correspondence between the Lean statements and the question,
and recomputed the main numbers with its own programs. A second independent
Claude Code agent then reviewed the whole paper: it checked the proofs and the
statements against the Lean file, recomputed the numbers with its own
programs, and checked the references. The programs, logs and Lean code are
available in version~1.3.0 of the repository~\cite{repo} (directories
\texttt{certificates/carenini/}, \texttt{lean/Research/} and \texttt{logs/}).
The author is not a professional mathematician.
%% AUTHOR-ROLE-SENTENCE: to be completed after the author has read the paper
The correctness of the results rests on the written proofs and, for the
results listed in Table~\ref{tab:lean}, also on the Lean formalisation; the
computations of Section~\ref{sec:comp} are additional cross-checks. AI systems
are not authors of this paper; the author takes full responsibility for its
content.

\begin{thebibliography}{99}
\bibitem{alon}
N.~Alon, Independent sets in regular graphs and sum-free subsets of finite
groups, \emph{Israel J.\ Math.} \textbf{73} (1991), 247--256.
\bibitem{br}
R.~R. Bahadur and R.~Ranga Rao, On deviations of the sample mean,
\emph{Ann.\ Math.\ Statist.} \textbf{31} (1960), 1015--1027,
\url{https://doi.org/10.1214/aoms/1177705674}.
\bibitem{car}
G.~Carenini, Counting almost independent sets in regular graphs,
arXiv:2609.28527v1, 22 September 2026.
\bibitem{dz}
A.~Dembo and O.~Zeitouni, \emph{Large Deviations Techniques and
Applications}, 2nd ed., Applications of Mathematics~38, Springer, New York,
1998.
\bibitem{repo}
H.~Dong, \emph{ai4math-results}, repository,
\url{https://github.com/Horace-Maxwell/ai4math-results}, version~1.3.0
(2026); archived at Zenodo under the concept DOI
\url{https://doi.org/10.5281/zenodo.22970254}.
\bibitem{fhm}
D.~R. Fulkerson, A.~J. Hoffman and M.~H. McAndrew, Some properties of graphs
with multiple edges, \emph{Canad.\ J.\ Math.} \textbf{17} (1965), 166--177.
\bibitem{hakimi}
S.~L. Hakimi, On realizability of a set of integers as degrees of the
vertices of a linear graph.~I, \emph{J.\ Soc.\ Indust.\ Appl.\ Math.}
\textbf{10} (1962), 496--506.
\bibitem{kahn}
J.~Kahn, An entropy approach to the hard-core model on bipartite graphs,
\emph{Combin.\ Probab.\ Comput.} \textbf{10} (2001), 219--237.
\bibitem{mantel}
W.~Mantel, Problem 28 (solution by H.~Gouwentak, W.~Mantel, J.~Teixeira de
Mattes, F.~Schuh and W.~A. Wythoff), \emph{Wiskundige Opgaven} \textbf{10}
(1907), 60--61.
\bibitem{mathlib}
The mathlib Community, The Lean mathematical library, in \emph{Proceedings of
the 9th ACM SIGPLAN International Conference on Certified Programs and Proofs}
(CPP '20), ACM, 2020, \url{https://doi.org/10.1145/3372885.3373824}.
\bibitem{oeis}
OEIS Foundation Inc., \emph{The On-Line Encyclopedia of Integer Sequences},
\url{https://oeis.org}, sequences A002851, A006820 and A006821, accessed 26
September 2026.
\bibitem{ross}
N.~Ross, Fundamentals of Stein's method, \emph{Probab.\ Surv.} \textbf{8}
(2011), 210--293, \url{https://doi.org/10.1214/11-PS182}; theorem numbers as
in arXiv:1109.1880v1.
\bibitem{sssz}
A.~Sah, M.~Sawhney, D.~Stoner and Y.~Zhao, A reverse Sidorenko inequality,
\emph{Invent.\ Math.} \textbf{221} (2020), 665--711,
\url{https://doi.org/10.1007/s00222-020-00956-9}; also arXiv:1809.09462v3.
\bibitem{seth}
C.~Seth, A tolerant independent set tester, in \emph{Proceedings of the 57th
Annual ACM Symposium on Theory of Computing} (STOC '25), 2025, 1031--1042,
\url{https://doi.org/10.1145/3717823.3718290}; also arXiv:2503.21441.
\bibitem{trurepo}
The Omega Institute, \texttt{trureturing} repository,
\url{https://github.com/the-omega-institute/trureturing}, accessed 26
September 2026.
\bibitem{zhao}
Y.~Zhao, The number of independent sets in a regular graph,
\emph{Combin.\ Probab.\ Comput.} \textbf{19} (2010), 315--320.
\end{thebibliography}
\end{document}
